\documentclass[11pt]{article}
\usepackage[utf8]{inputenc}
\usepackage[margin=1in]{geometry}
\usepackage{amsmath}
\usepackage{graphicx}
\usepackage{booktabs}
\usepackage{hyperref}
\usepackage[round]{natbib}

\title{A Translational MRI Approach to Validate Acute Axonal Damage Detection Using AxCaliber Modeling}
\author{Author A$^{1,2}$, Author B$^{1}$, Author C$^{2}$, Author D$^{3}$ \\
\small $^{1}$Department of Neurology, Tel Aviv University, Tel Aviv, Israel \\
\small $^{2}$Center for Advanced MRI, Tel Aviv Sourasky Medical Center, Israel \\
\small $^{3}$Department of Neurobiology, Weizmann Institute of Science, Rehovot, Israel}
\date{}

\begin{document}
\maketitle

\begin{abstract}
Axonal degeneration is a central pathological feature driving irreversible disability in neurodegenerative diseases such as multiple sclerosis (MS), yet current noninvasive detection methods lack specificity for axonal damage. Here we present a translational MRI approach to validate AxCaliber---a multicompartment diffusion MRI model estimating axonal diameter distributions---as a sensitive and specific biomarker for acute axonal damage. In the preclinical arm, unilateral stereotaxic injection of ibotenic acid into the rat hippocampus ($n = 9$) produced selective axonal degeneration in the fimbria while preserving myelination. AxCaliber revealed significantly increased mean axonal diameter in the injected versus control fimbria (paired $t$-test, $p = 0.003$), with histological validation showing significant correlation between neurofilament fluorescence intensity and MRI-estimated axonal diameter ($r = 0.54$, $p = 0.029$). Critically, myelin basic protein (MBP) staining showed no significant differences, confirming axonal specificity. In the clinical arm, whole-brain AxCaliber mapping in 10 relapsing-remitting MS patients and 6 healthy controls revealed widespread, diffuse increases in mean axonal diameter across most normal-appearing white matter (NAWM) regions, and significantly elevated normalized index of asymmetry in MS lesions. A negative association between axonal diameter and disease duration suggests early axonal swelling followed by progressive axonal loss. These findings validate AxCaliber as a sensitive and specific tool for detecting acute axonal damage and reveal diffuse axonal pathology in MS NAWM.
\end{abstract}

\section{Introduction}

Axonal degeneration is increasingly recognized as the principal pathological substrate of irreversible neurological disability in multiple sclerosis and other neurodegenerative diseases \citep{trapp1998axonal, dutta2011mechanisms}. While demyelination has traditionally received the most attention in MS research, evidence suggests that axonal damage occurs independently of and often precedes myelin loss, and that the extent of axonal degeneration is the strongest correlate of permanent disability \citep{deStephano2001axonal, kutzelnigg2005cortical}.

Current clinical MRI techniques, including conventional diffusion tensor imaging (DTI), lack specificity for axonal damage because changes in fractional anisotropy and mean diffusivity can reflect a mixture of microstructural alterations including demyelination, inflammation, edema, and axonal pathology \citep{alexander2007diffusion}. This limitation has motivated the development of advanced diffusion MRI models that aim to separate contributions from different tissue compartments.

AxCaliber is a multicompartment diffusion MRI model that estimates axonal diameter distributions by separating intra-axonal from extra-axonal water compartments using diffusion-weighted signals acquired at multiple b-values and diffusion times \citep{assaf2008axcaliber}. While AxCaliber has been applied to measure axonal diameter in the corpus callosum of MS patients, these studies lacked independent histological validation. Furthermore, the original AxCaliber formulation was limited to single-fiber-orientation voxels (approximately 30\% of brain white matter), since crossing fibers confound the model. Recent extensions have overcome this limitation by incorporating fiber orientation distributions, enabling whole-brain axonal diameter mapping \citep{barazany2009vivo}.

A critical gap in the literature is the lack of controlled preclinical validation specifically targeting isolated axonal damage without confounding demyelination. Here we address this gap using a rat model of selective axonal degeneration induced by ibotenic acid, an NMDA receptor agonist that causes excitotoxic death of neuronal cell bodies in the hippocampus, leading to Wallerian-like degeneration of their axons in the fimbria while initially preserving myelin. We combine this preclinical validation with a clinical study applying whole-brain AxCaliber to MS patients, providing a translational framework for axonal damage detection.

\section{Methods}

\subsection{Study Design}

The study comprised two arms: a preclinical arm using a rat model of acute axonal damage with histological validation, and a clinical arm applying AxCaliber to MS patients and healthy controls (Table~\ref{tab:design}).

\begin{table}[htbp]
\centering
\caption{Study design overview.}
\label{tab:design}
\begin{tabular}{llll}
\toprule
Parameter & Preclinical Arm & Clinical Arm \\
\midrule
Subjects & $n = 9$ rats & 10 MS, 6 controls \\
Model/Pathology & Ibotenic acid injection & Relapsing-remitting MS \\
Target structure & Fimbria & Whole-brain WM \\
MRI protocol & AxCaliber & AxCaliber (clinical) \\
Control & Contralateral saline & Age-matched healthy \\
Validation & NeuN, neurofilament, MBP & N/A \\
\bottomrule
\end{tabular}
\end{table}

\subsection{Preclinical Model}

Nine rats received unilateral stereotaxic injection of ibotenic acid (2.5 $\mu$g/$\mu$L, 1 $\mu$L) into the right dorsal hippocampus at coordinates bregma $-3.8$ mm, superior-inferior 3.0 mm, 2 mm from midline. The contralateral hemisphere received saline injection as an internal control. Fourteen days post-surgery, rats underwent the AxCaliber MRI protocol followed by perfusion and immunohistochemical analysis.

Immunohistochemistry employed three stains: NeuN for neuronal cell bodies (hippocampus), neurofilament for axonal integrity (fimbria), and myelin basic protein (MBP) for myelination (fimbria). Fluorescence intensity was quantified and compared between injected and control hemispheres using paired $t$-tests.

\subsection{Clinical Cohort}

Ten relapsing-remitting MS patients (age 27--59, 7 female) and 6 age-matched healthy controls (age 23--53, 3 female) underwent whole-brain AxCaliber MRI. Age matching was assessed by Mann-Whitney test ($p = 0.093$). Tract-based spatial statistics (TBSS) compared NAWM axonal diameter between groups. For lesion analysis, a normalized index of asymmetry (NIA) was computed between lesion voxels and contralateral NAWM in MS patients, and compared to the same anatomical regions in healthy controls (Table~\ref{tab:cohort}).

\begin{table}[htbp]
\centering
\caption{Clinical cohort demographics.}
\label{tab:cohort}
\begin{tabular}{llll}
\toprule
Characteristic & MS Patients & Healthy Controls & $p$-value \\
\midrule
$N$ & 10 & 6 & -- \\
Age range & 27--59 years & 23--53 years & 0.093 \\
Sex (female) & 7/10 & 3/6 & -- \\
\bottomrule
\end{tabular}
\end{table}

\subsection{AxCaliber MRI Protocol}

AxCaliber is a multicompartment diffusion model that estimates mean axonal diameter from diffusion-weighted signals acquired at multiple b-values and diffusion times. The model separates intra-axonal from extra-axonal water compartments, with the intra-axonal compartment modeled as restricted diffusion within cylinders of varying diameter. The whole-brain extension incorporates fiber orientation distributions to handle crossing-fiber voxels, overcoming the original single-fiber limitation.

\subsection{Statistical Analysis}

Preclinical analyses used paired $t$-tests for hemisphere comparisons and repeated-measures ANOVA for along-tract analysis with factors of injection (injected vs. control) and position (100 steps along the antero-posterior axis). Pearson correlation assessed the relationship between neurofilament fluorescence intensity and MRI-estimated axonal diameter. Clinical analyses used nonparametric permutation testing for TBSS group comparisons and correlation/regression for disease duration associations (Table~\ref{tab:stats}).

\begin{table}[htbp]
\centering
\caption{Statistical methods summary.}
\label{tab:stats}
\begin{tabular}{lll}
\toprule
Analysis & Test Used & Context \\
\midrule
Injected vs. control diameter & Paired $t$-test & Preclinical fimbria \\
Along-tract analysis & RM-ANOVA + post-hoc & Position $\times$ injection \\
NeuN intensity & Paired $t$-test & Hippocampal neuronal loss \\
Neurofilament intensity & Paired $t$-test & Fimbria axonal damage \\
MRI-histology correlation & Pearson correlation & Neurofilament vs. diameter \\
Age comparison & Mann-Whitney & Cohort matching \\
TBSS group comparison & Permutation testing & Whole-brain NAWM \\
Disease duration & Correlation/regression & TBSS and ROI-based \\
\bottomrule
\end{tabular}
\end{table}

\section{Results}

\subsection{Ibotenic Acid Induces Selective Axonal Damage}

NeuN staining confirmed significant neuronal loss in the injected hippocampus compared to the control hemisphere (paired $t$-test, $p = 0.026$), validating the excitotoxic lesion. Neurofilament staining in the fimbria showed significantly higher fluorescence intensity on the injected side ($p = 0.047$), consistent with axonal swelling and damage. Critically, MBP staining revealed no significant differences between injected and control fimbria, confirming that myelin was preserved at 14 days post-injection despite clear axonal pathology (Table~\ref{tab:histology}).

\begin{table}[htbp]
\centering
\caption{Immunohistochemistry results comparing injected versus control hemispheres.}
\label{tab:histology}
\begin{tabular}{lllll}
\toprule
Stain & Target & Finding & $p$-value & Test \\
\midrule
NeuN & Neuronal somas & Lower (neuronal loss) & 0.026 & Paired $t$-test \\
Neurofilament & Axonal integrity & Higher (axonal damage) & 0.047 & Paired $t$-test \\
MBP & Myelin & No significant difference & NS & Paired $t$-test \\
\bottomrule
\end{tabular}
\end{table}

\subsection{AxCaliber Detects Axonal Diameter Changes}

AxCaliber MRI revealed significantly increased mean axonal diameter in the ibotenic acid-injected fimbria compared to the saline-injected control (paired $t$-test, $p = 0.003$). Along-tract analysis using repeated-measures ANOVA confirmed significant main effects of injection ($F = 19.5$, $df = 1,7$, $p = 0.003$), position along tract ($F = 219.5$, $df = 98,686$, $p < 0.001$), and their interaction ($F = 11.2$, $df = 98,686$, $p = 0.012$). The diameter differences were primarily localized posterior to the injection site, consistent with the anatomical trajectory of degenerating axons.

\subsection{MRI-Histology Correlation Validates AxCaliber}

A significant positive correlation was observed between neurofilament fluorescence intensity and AxCaliber-estimated axonal diameter across both injected and control fimbria ($r = 0.54$, $p = 0.029$). This correlation provides direct histological validation that MRI-based axonal diameter estimation reflects underlying axonal morphological changes, supporting the use of AxCaliber as a biomarker for axonal pathology.

\subsection{Diffuse NAWM Axonal Diameter Increases in MS}

TBSS analysis revealed widespread, mostly symmetrical increases in mean axonal diameter in MS NAWM compared to healthy controls ($p < 0.05$, corrected). Affected white matter tracts included the corpus callosum, internal capsule, corona radiata, thalamic radiation, inferior longitudinal fasciculus, cingulum, fornix, superior longitudinal fasciculus, inferior fronto-occipital fasciculus, uncinate fasciculus, and tapetum. No significant decreases were observed.

\subsection{Elevated Axonal Diameter in MS Lesions}

Analysis of the normalized index of asymmetry between MS lesions and contralateral NAWM revealed significantly elevated NIA in MS patients compared to the same anatomical regions in healthy controls. This result confirms that axonal pathology is more pronounced within lesions, consistent with the expected pattern of focal inflammatory demyelinating damage causing additional axonal injury.

\subsection{Negative Association with Disease Duration}

Both TBSS and ROI-based analyses revealed a negative association between axonal diameter and disease duration. This pattern suggests that early disease stages are characterized by axonal swelling (varicosities and spheroids that enlarge the measured diameter and impair axonal transport), while progressive axonal loss at later stages removes the swollen axons, reducing the mean diameter. Additionally, smaller axons may be preferentially lost earlier, initially biasing the mean upward before overall axonal density decreases.

\section{Discussion}

\subsection{Validation of AxCaliber for Axonal Damage Detection}

Our translational approach provides the first controlled validation of AxCaliber sensitivity and specificity for acute axonal damage. The preclinical model demonstrates three key findings: (1) AxCaliber detects increased axonal diameter in the presence of axonal damage ($p = 0.003$), (2) MRI estimates correlate with histological measures of axonal integrity ($r = 0.54$), and (3) the changes are specific to axonal damage, as myelin is preserved. These findings establish AxCaliber as a validated biomarker with known sensitivity and specificity characteristics.

\subsection{Diffuse Axonal Pathology in MS NAWM}

The finding of widespread increased axonal diameter in MS NAWM has important implications for understanding MS pathogenesis. The conventional view that NAWM is largely unaffected is challenged by histological studies showing subclinical axonal alterations even in regions without visible myelin or inflammatory abnormalities \citep{kutzelnigg2005cortical}. Our AxCaliber results provide in vivo corroboration of these histological findings, suggesting that diffuse axonal pathology is a pervasive feature of relapsing-remitting MS that extends well beyond visible lesion boundaries.

\subsection{Temporal Dynamics of Axonal Damage}

The negative association between axonal diameter and disease duration suggests a biphasic process: initial axonal swelling (detectable as increased diameter) followed by progressive loss of damaged axons (detectable as decreased diameter or decreased axonal density). This temporal framework has implications for clinical monitoring and treatment timing---interventions targeting axonal protection may be most effective during the early swelling phase before irreversible axonal loss occurs.

\subsection{Advantages Over Conventional DTI}

AxCaliber offers a fundamental advantage over DTI by providing a specific measure of the intra-axonal compartment rather than a composite measure affected by multiple tissue properties \citep{alexander2007diffusion}. This specificity is critical for distinguishing axonal damage from demyelination, inflammation, and edema---a distinction that has direct implications for treatment selection and prognostication.

\subsection{Limitations}

Several limitations should be noted. The preclinical model uses an acute excitotoxic lesion that may not fully recapitulate the chronic, multifocal axonal damage in MS. The clinical sample size is modest ($n = 10$ MS, $n = 6$ controls), limiting statistical power for subgroup analyses. The AxCaliber model makes assumptions about axonal geometry (cylindrical shape, gamma-distributed diameters) that may not perfectly capture the complexity of real tissue. Fixation artifacts in histological tissue prevent perfect quantitative comparison between MRI and histology.

\section{Conclusion}

We have demonstrated that AxCaliber MRI provides a sensitive and specific biomarker for acute axonal damage through a translational validation approach combining controlled preclinical experiments with clinical application. The significant correlation between AxCaliber estimates and histological measures of axonal integrity, combined with the absence of myelin changes, establishes the specificity of this approach for axonal pathology. Application to MS reveals diffuse axonal damage throughout NAWM, challenging the conventional focus on focal lesions and suggesting that axonal protection strategies should target the entire white matter. The temporal relationship between axonal diameter and disease duration provides a framework for understanding the progression of axonal pathology in MS.

\bibliographystyle{plainnat}
\bibliography{references}

\end{document}
